% §11.3 RP, BPP, and Bernoulli Type
%
% Monograph register + the RP/coRP axis fix. The draft swapped the axis labels:
% RP has eps = 0 (FPR-ZERO axis, certain accept, the perfect-hash-filter shape),
% coRP has omega = 0 (FNR-ZERO axis, certain reject, the BLOOM-filter shape).
% Miller-Rabin puts Composites in RP (certain composite-witness) and Primes in
% coRP, consistent with §11.1/§11.2/§11.4 (which correctly frame MR-primality as
% omega = 0 / FNR-zero). The table VALUES, the Schwartz-Zippel example, the PPV
% discussion, and both amplification theorems were already correct. Cross-refs
% def:bernoulli-bool, thm:composition, ax:element-indep, def:perfect-filter-line,
% sec:perfect-filter-line, def:ppv, def:npv, sec:one-sided-two-sided,
% def:bloom-filter, ch:perfect-hash-filters, thm:miller-rabin-formal,
% def:monte-carlo, sec:sample-complexity, thm:chernoff. DEFINED: sec:rp-bpp,
% def:rp, def:corp, def:zpp, def:bpp, thm:rp-amplification, thm:bpp-amplification.

\section{RP, BPP, and Bernoulli Type}
\label{sec:rp-bpp}

The Monte Carlo decision problems (\Cref{def:monte-carlo}) carve into the
classes RP, coRP, ZPP, and BPP by where the error is allowed to land. Each is a
location on the chapter-6 rate plane: RP and coRP on the two axes of the perfect
filter line, ZPP at the origin, BPP in the interior. The correspondence is not
an analogy; the complexity-class definitions and the chapter-6 classification are
the same axes from two angles.

\subsection*{RP: certain accept, the FPR-zero axis}

\begin{definition}[Randomized Polynomial Time, RP]
\label{def:rp}
A language $L$ is in \emph{RP} when a polynomial-time randomized $A$ satisfies
$x \in L \Rightarrow \Prob{A(x) = 1} \ge 1/2$ and $x \notin L \Rightarrow
\Prob{A(x) = 1} = 0$, the probability over $A$'s coins.
\end{definition}

Read $A(x) = 1$ as the report \textsc{in-}$L$. The second line says the algorithm
never accepts a non-member, $\fprate = \Prob{A(x) = 1 \mid x \notin L} = 0$; the
first says it accepts a member with probability $\ge 1/2$, $\fnrate = \Prob{A(x)
= 0 \mid x \in L} \le 1/2$. An RP algorithm is a $\Bernoulli{\BoolType}$ instance
with $\fprate = 0$, $\fnrate \le 1/2$: it sits on the \emph{FPR-zero} axis of the
perfect filter line (\Cref{def:perfect-filter-line}). An accept is certain, a
witness seen is a witness believed, and the only failure is missing a witness
that exists. This is the certain-accept shape, the decision-algorithm kin of the
perfect hash filter of \Cref{ch:perfect-hash-filters}, not the Bloom filter,
whose certainty is on the reject side.

The canonical RP problem is \textsc{Non-Zero Polynomial}: given a circuit for a
polynomial $f$ over a finite field, decide whether $f \not\equiv 0$. By
Schwartz--Zippel, a random assignment makes a non-zero $f$ evaluate non-zero
with probability $\ge 1 - d/|F|$ (degree $d$, field size $|F|$), while a zero
$f$ always evaluates zero. Reporting \textsc{non-zero} on a non-zero evaluation
is the RP shape: silent on zero polynomials ($\fprate = 0$), occasionally
missing a non-zero one ($\fnrate \le d/|F|$).

\subsection*{coRP and ZPP: the mirror axis and the origin}

\begin{definition}[coRP and ZPP]
\label{def:corp}
$L \in \emph{coRP}$ when $\bar L \in \mathrm{RP}$: equivalently $x \in L
\Rightarrow \Prob{A(x) = 1} = 1$ and $x \notin L \Rightarrow \Prob{A(x) = 1} \le
1/2$. \label{def:zpp} $L \in \emph{ZPP}$ when $L \in \mathrm{RP} \cap
\mathrm{coRP}$.
\end{definition}

A coRP algorithm has $\fnrate = 0$, $\fprate \le 1/2$: it sits on the
\emph{FNR-zero} axis, the certain-\emph{reject} shape, the decision-algorithm kin
of the Bloom filter (\Cref{def:bloom-filter}), whose $\bot$ report is certain.
Miller-Rabin (\Cref{thm:miller-rabin-formal}) is the canonical example: a
composite witness certifies compositeness, so on the language \textsc{Primes} a
reject is certain ($\fnrate = 0$, a prime is never rejected) while a composite is
occasionally accepted as prime ($\fprate \le 1/4$). Thus $\textsc{Primes} \in
\mathrm{coRP}$, and dually $\textsc{Composites} \in \mathrm{RP}$ via the same
algorithm, where instead a composite-witness is a certain \emph{accept} of
compositeness. (This was the textbook RP/coRP example until the deterministic
AKS test of 2002.) ZPP, at $\fprate = \fnrate = 0$, is the origin: exact
decisions paid for in the runtime tail, equivalently the expected-polynomial-time
Las Vegas class.

\subsection*{BPP: the interior}

\begin{definition}[Bounded-error Probabilistic Polynomial Time, BPP]
\label{def:bpp}
$L \in \emph{BPP}$ when a polynomial-time randomized $A$ satisfies $x \in L
\Rightarrow \Prob{A(x) = 1} \ge 2/3$ and $x \notin L \Rightarrow \Prob{A(x) = 1}
\le 1/3$.
\end{definition}

Read as a channel, $\fnrate, \fprate \le 1/3$: BPP sits in the interior of the
rate plane, off both axes. The chapter-7 vocabulary applies directly. At
prevalence $\pi = \Prob{x \in L}$ the positive predictive value
(\Cref{def:ppv}) is
\[
  \PPV = \frac{\pi(1 - \fnrate)}{\pi(1 - \fnrate) + (1 - \pi)\fprate}
       \ge \frac{\pi \cdot 2/3}{\pi \cdot 2/3 + (1 - \pi)\cdot 1/3}
       = \frac{2\pi}{\pi + 1},
\]
which at $\pi = 0.01$ is only $\approx 0.0196$: despite $2/3$-versus-$1/3$
per-class accuracy, $98\%$ of accept reports are false at low prevalence, the
same prevalence-dependence \Cref{def:npv} and chapter~7 record. The $2/3$-versus-
$1/3$ gap is conventional but load-bearing: it forces a signal in both
directions so amplification can drive both errors down, where the gap-closed
class PP cannot.

\subsection*{Amplification}

\begin{theorem}[RP amplification]
\label{thm:rp-amplification}
For an RP algorithm with $\fnrate \le 1/2$, running $k$ independent times and
accepting if any run accepts gives an RP algorithm with $\fnrate \le 2^{-k}$,
$\fprate = 0$.
\end{theorem}

This is the composition theorem (\Cref{thm:composition}) with the OR aggregator:
the certain-accept property ($\fprate = 0$) is preserved (any accept is true),
and $\fnrate$ at the OR is the product of per-run rates by independence
(\Cref{ax:element-indep}). It is the FPR-zero mirror of the $k$-round
Miller-Rabin bound, which drives the coRP rate $\fprate$ down via AND; coRP
amplifies dually. Bringing $\fnrate$ to $2^{-30}$ costs $30$ runs.

\begin{theorem}[BPP amplification]
\label{thm:bpp-amplification}
For a BPP algorithm with $\fnrate, \fprate \le 1/3$, the majority of $k$
independent runs gives $\fnrate, \fprate \le \exp(-\Omega(k))$.
\end{theorem}

The exponent is no longer the per-run rate but the Chernoff exponent for a
majority of independent draws; the proof, via the Chernoff bound
(\Cref{thm:chernoff}), is in \Cref{sec:sample-complexity}. BPP error decays
exponentially in the repetition count even though no single coin sequence drives
either error to zero.

\subsection*{The correspondence at scale}

\begin{center}
\begin{tabular}{lcc}
\hline
\textbf{Class} & $\fprate$ bound & $\fnrate$ bound \\
\hline
RP   & $= 0$       & $\le 1/2$ \\
coRP & $\le 1/2$   & $= 0$ \\
ZPP  & $= 0$       & $= 0$ \\
BPP  & $\le 1/3$   & $\le 1/3$ \\
\hline
\end{tabular}
\end{center}
The first three rows trace the perfect filter line of
\Cref{def:perfect-filter-line}: RP the FPR-zero axis, coRP the FNR-zero axis,
ZPP their intersection at the origin. BPP is the interior below $1/3$. The
chapter-6 classification of approximate sets by axis and the classical
randomized-complexity taxonomy are one classification applied to two families,
arbitrary $\Bernoulli{\BoolType}$ instances there, polynomially constrained
decision algorithms here. The one-sided-versus-two-sided distinction of
\Cref{sec:one-sided-two-sided} is what separates RP and coRP from BPP, and
amplification is the constraint that propagates: any instance with per-trial
error bounded away from $1/2$ on both inputs is driven to negligible error by
repetition and the right aggregator (OR or AND on an axis, majority in the
interior). Chapters~12 and~13 lift these classes along with the rest of the
framework, so the relational analog of BPP amplifies for the same reason.
